\documentclass[a4paper,12pt,twoside]{memoir}

% Castellano
\usepackage[spanish,es-tabla]{babel}
\selectlanguage{spanish}
\usepackage[utf8]{inputenc}
\usepackage[T1]{fontenc}
\usepackage{lmodern} % Scalable font
\usepackage{microtype}
\usepackage{placeins}

\RequirePackage{booktabs}
\RequirePackage[table]{xcolor}
\RequirePackage{xtab}
\RequirePackage{multirow}

% Links
\PassOptionsToPackage{hyphens}{url}\usepackage[colorlinks]{hyperref}
\hypersetup{
	allcolors = {red}
}

% Ecuaciones
\usepackage{amsmath}

% Rutas de fichero / paquete
\newcommand{\ruta}[1]{{\sffamily #1}}

% Párrafos
\nonzeroparskip

% Huérfanas y viudas
\widowpenalty100000
\clubpenalty100000

% Imágenes

% Comando para insertar una imagen en un lugar concreto.
% Los parámetros son:
% 1 --> Ruta absoluta/relativa de la figura
% 2 --> Texto a pie de figura
% 3 --> Tamaño en tanto por uno relativo al ancho de página
\usepackage{graphicx}
\newcommand{\imagen}[3]{
	\begin{figure}[!h]
		\centering
		\includegraphics[width=#3\textwidth]{#1}
		\caption{#2}\label{fig:#1}
	\end{figure}
	\FloatBarrier
}

% Comando para insertar una imagen sin posición.
% Los parámetros son:
% 1 --> Ruta absoluta/relativa de la figura
% 2 --> Texto a pie de figura
% 3 --> Tamaño en tanto por uno relativo al ancho de página
\newcommand{\imagenflotante}[3]{
	\begin{figure}
		\centering
		\includegraphics[width=#3\textwidth]{#1}
		\caption{#2}\label{fig:#1}
	\end{figure}
}

% El comando \figura nos permite insertar figuras comodamente, y utilizando
% siempre el mismo formato. Los parametros son:
% 1 --> Porcentaje del ancho de página que ocupará la figura (de 0 a 1)
% 2 --> Fichero de la imagen
% 3 --> Texto a pie de imagen
% 4 --> Etiqueta (label) para referencias
% 5 --> Opciones que queramos pasarle al \includegraphics
% 6 --> Opciones de posicionamiento a pasarle a \begin{figure}
\newcommand{\figuraConPosicion}[6]{%
  \setlength{\anchoFloat}{#1\textwidth}%
  \addtolength{\anchoFloat}{-4\fboxsep}%
  \setlength{\anchoFigura}{\anchoFloat}%
  \begin{figure}[#6]
    \begin{center}%
      \Ovalbox{%
        \begin{minipage}{\anchoFloat}%
          \begin{center}%
            \includegraphics[width=\anchoFigura,#5]{#2}%
            \caption{#3}%
            \label{#4}%
          \end{center}%
        \end{minipage}
      }%
    \end{center}%
  \end{figure}%
}

%
% Comando para incluir imágenes en formato apaisado (sin marco).
\newcommand{\figuraApaisadaSinMarco}[5]{%
  \begin{figure}%
    \begin{center}%
    \includegraphics[angle=90,height=#1\textheight,#5]{#2}%
    \caption{#3}%
    \label{#4}%
    \end{center}%
  \end{figure}%
}
% Para las tablas
\newcommand{\otoprule}{\midrule [\heavyrulewidth]}
%
% Nuevo comando para tablas pequeñas (menos de una página).
\newcommand{\tablaSmall}[5]{%
 \begin{table}
  \begin{center}
   \rowcolors {2}{gray!35}{}
   \begin{tabular}{#2}
    \toprule
    #4
    \otoprule
    #5
    \bottomrule
   \end{tabular}
   \caption{#1}
   \label{tabla:#3}
  \end{center}
 \end{table}
}

%
% Nuevo comando para tablas pequeñas (menos de una página).
\newcommand{\tablaSmallSinColores}[5]{%
 \begin{table}[H]
  \begin{center}
   \begin{tabular}{#2}
    \toprule
    #4
    \otoprule
    #5
    \bottomrule
   \end{tabular}
   \caption{#1}
   \label{tabla:#3}
  \end{center}
 \end{table}
}

\newcommand{\tablaApaisadaSmall}[5]{%
\begin{landscape}
  \begin{table}
   \begin{center}
    \rowcolors {2}{gray!35}{}
    \begin{tabular}{#2}
     \toprule
     #4
     \otoprule
     #5
     \bottomrule
    \end{tabular}
    \caption{#1}
    \label{tabla:#3}
   \end{center}
  \end{table}
\end{landscape}
}

%
% Nuevo comando para tablas grandes con cabecera y filas alternas coloreadas en gris.
\newcommand{\tabla}[6]{%
  \begin{center}
    \tablefirsthead{
      \toprule
      #5
      \otoprule
    }
    \tablehead{
      \multicolumn{#3}{l}{\small\sl continúa desde la página anterior}\\
      \toprule
      #5
      \otoprule
    }
    \tabletail{
      \hline
      \multicolumn{#3}{r}{\small\sl continúa en la página siguiente}\\
    }
    \tablelasttail{
      \hline
    }
    \bottomcaption{#1}
    \rowcolors {2}{gray!35}{}
    \begin{xtabular}{#2}
      #6
      \bottomrule
    \end{xtabular}
    \label{tabla:#4}
  \end{center}
}

%
% Nuevo comando para tablas grandes con cabecera.
\newcommand{\tablaSinColores}[6]{%
  \begin{center}
    \tablefirsthead{
      \toprule
      #5
      \otoprule
    }
    \tablehead{
      \multicolumn{#3}{l}{\small\sl continúa desde la página anterior}\\
      \toprule
      #5
      \otoprule
    }
    \tabletail{
      \hline
      \multicolumn{#3}{r}{\small\sl continúa en la página siguiente}\\
    }
    \tablelasttail{
      \hline
    }
    \bottomcaption{#1}
    \begin{xtabular}{#2}
      #6
      \bottomrule
    \end{xtabular}
    \label{tabla:#4}
  \end{center}
}

%
% Nuevo comando para tablas grandes sin cabecera.
\newcommand{\tablaSinCabecera}[5]{%
  \begin{center}
    \tablefirsthead{
      \toprule
    }
    \tablehead{
      \multicolumn{#3}{l}{\small\sl continúa desde la página anterior}\\
      \hline
    }
    \tabletail{
      \hline
      \multicolumn{#3}{r}{\small\sl continúa en la página siguiente}\\
    }
    \tablelasttail{
      \hline
    }
    \bottomcaption{#1}
  \begin{xtabular}{#2}
    #5
   \bottomrule
  \end{xtabular}
  \label{tabla:#4}
  \end{center}
}



\definecolor{cgoLight}{HTML}{EEEEEE}
\definecolor{cgoExtralight}{HTML}{FFFFFF}

%
% Nuevo comando para tablas grandes sin cabecera.
\newcommand{\tablaSinCabeceraConBandas}[5]{%
  \begin{center}
    \tablefirsthead{
      \toprule
    }
    \tablehead{
      \multicolumn{#3}{l}{\small\sl continúa desde la página anterior}\\
      \hline
    }
    \tabletail{
      \hline
      \multicolumn{#3}{r}{\small\sl continúa en la página siguiente}\\
    }
    \tablelasttail{
      \hline
    }
    \bottomcaption{#1}
    \rowcolors[]{1}{cgoExtralight}{cgoLight}

  \begin{xtabular}{#2}
    #5
   \bottomrule
  \end{xtabular}
  \label{tabla:#4}
  \end{center}
}



\graphicspath{ {./img/} }

% Capítulos
\chapterstyle{bianchi}
\newcommand{\capitulo}[2]{
	\setcounter{chapter}{#1}
	\setcounter{section}{0}
	\setcounter{figure}{0}
	\setcounter{table}{0}
	\chapter*{#2}
	\addcontentsline{toc}{chapter}{#2}
	\markboth{#2}{#2}
}

% Apéndices
\renewcommand{\appendixname}{Apéndice}
\renewcommand*\cftappendixname{\appendixname}

\newcommand{\apendice}[1]{
	%\renewcommand{\thechapter}{A}
	\chapter{#1}
}

\renewcommand*\cftappendixname{\appendixname\ }

% Formato de portada
\makeatletter
\usepackage{xcolor}
\newcommand{\tutor}[1]{\def\@tutor{#1}}
\newcommand{\course}[1]{\def\@course{#1}}
\definecolor{cpardoBox}{HTML}{E6E6FF}
\def\maketitle{
  \null
  \thispagestyle{empty}
  % Cabecera ----------------
\noindent\includegraphics[width=\textwidth]{cabecera}\vspace{1cm}%
  \vfill
  % Título proyecto y escudo informática ----------------
  \colorbox{cpardoBox}{%
    \begin{minipage}{.8\textwidth}
      \vspace{.5cm}\Large
      \begin{center}
      \textbf{TFG del Grado en Ingeniería Informática}\vspace{.6cm}\\
      \textbf{\LARGE\@title{}}
      \end{center}
      \vspace{.2cm}
    \end{minipage}

  }%
  \hfill\begin{minipage}{.20\textwidth}
    \includegraphics[width=\textwidth]{escudoInfor}
  \end{minipage}
  \vfill
  % Datos de alumno, curso y tutores ------------------
  \begin{center}%
  {%
    \noindent\LARGE
    Presentado por \@author{}\\ 
    en Universidad de Burgos --- \@date{}\\
    Tutor: \@tutor{}\\
  }%
  \end{center}%
  \null
  \cleardoublepage
  }
\makeatother

\newcommand{\nombre}{Nombre del alumno} %%% cambio de comando

% Datos de portada
\title{título del TFG}
\author{\nombre}
\tutor{nombre tutor}
\date{\today}

\begin{document}

\maketitle


\newpage\null\thispagestyle{empty}\newpage


%%%%%%%%%%%%%%%%%%%%%%%%%%%%%%%%%%%%%%%%%%%%%%%%%%%%%%%%%%%%%%%%%%%%%%%%%%%%%%%%%%%%%%%%
\thispagestyle{empty}


\noindent\includegraphics[width=\textwidth]{cabecera}\vspace{1cm}

\noindent D. nombre tutor, profesor del departamento de nombre departamento, área de nombre área.

\noindent Expone:

\noindent Que el alumno D. \nombre, con DNI dni, ha realizado el Trabajo final de Grado en Ingeniería Informática titulado título de TFG. 

\noindent Y que dicho trabajo ha sido realizado por el alumno bajo la dirección del que suscribe, en virtud de lo cual se autoriza su presentación y defensa.

\begin{center} %\large
En Burgos, {\large \today}
\end{center}

\vfill\vfill\vfill

% Author and supervisor
\begin{minipage}{0.45\textwidth}
\begin{flushleft} %\large
Vº. Bº. del Tutor:\\[2cm]
D. nombre tutor
\end{flushleft}
\end{minipage}
\hfill
\begin{minipage}{0.45\textwidth}
\begin{flushleft} %\large
Vº. Bº. del co-tutor:\\[2cm]
D. nombre co-tutor
\end{flushleft}
\end{minipage}
\hfill

\vfill

% para casos con solo un tutor comentar lo anterior
% y descomentar lo siguiente
%Vº. Bº. del Tutor:\\[2cm]
%D. nombre tutor


\newpage\null\thispagestyle{empty}\newpage




\frontmatter

% Abstract en castellano
\renewcommand*\abstractname{Resumen}
\begin{abstract}
En este primer apartado se hace una \textbf{breve} presentación del tema que se aborda en el proyecto.
\end{abstract}

\renewcommand*\abstractname{Descriptores}
\begin{abstract}
Palabras separadas por comas que identifiquen el contenido del proyecto Ej: servidor web, buscador de vuelos, android \ldots
\end{abstract}

\clearpage

% Abstract en inglés
\renewcommand*\abstractname{Abstract}
\begin{abstract}
A \textbf{brief} presentation of the topic addressed in the project.
\end{abstract}

\renewcommand*\abstractname{Keywords}
\begin{abstract}
keywords separated by commas.
\end{abstract}

\clearpage

% Indices
\tableofcontents

\clearpage

\listoffigures

\clearpage

\listoftables
\clearpage

\mainmatter
\capitulo{1}{Introducción}
El cuidado de las tierras de cultivo vitícolas ha supuesto a menudo a lo largo 
de la historia diferentes retos desde el comienzo de su domesticación entorno al 
3500-3100 a.C., continuando en la época del Antiguo Egipto, lugar en el que se cree se comenzó a popularizar la ingesta
del producto resultante de la labranza y cosecha de los cultivos: el vino \cite{book:piqueras2014}.

En la península ibérica el interés por este proceso y su resultante comenzó en el 
periodo fenicio cerca del 1100 a.C., expandiéndose por todo el territorio peninsular
gracias al comercio, lo que dio como resultado que los procesos enológicos fueran ya bien
conocidos en los siglos IV-III a.C. \cite{book:piqueras2014, MAPA:Historia}.

Entre los principales retos a los que se han enfrentado los viticultores en la península ibérica se 
encuentran numerosas plagas de insectos, ácaros, nematodos, diferentes vertebrados, moluscos, bacterias, etc \cite{book:mundi2004cap1}.
Siendo la más devastadora la causada por la histórica plaga de filoxera (\textit{Phylloxera vastatrix}) a 
finales del siglo XIX y comienzos del XX, que obligó a una reestructuración prácticamente completa de la viticultura española \cite{book:mundi2004cap3}.

Sin embargo, el principal factor influyente en el cultivo de la vid es el clima, siendo este el causante de sus alteraciones, fisiopatías y
la presencia de casi todas las plagas, de manera que estas se encuentran relacionadas con diferentes 
factores climáticos como la temperatura, las precipitaciones, la humedad y la humectación. 
De esta forma, por ejemplo, las altas lluvias o humedades y/o bajas temperaturas pueden producir corrimiento del racimo (escasez anormal de bayas en los racimos) o 
las altas temperaturas pueden generar desequilibrios hídricos que afecten a los injertos y a aquellas cepas infectadas por hongos de la madera \cite{book:mundi2004cap1}.

\clearpage

Por otro lado, los datos climatológicos históricos indican que la temperatura global ha aumentado 1,1ºC desde el 
inicio de la Revolución Industrial hasta nuestros días, siendo este desequilibrio causado en los últimos años
por la actividad humana, produciéndose a velocidad tal que los ecosistemas no tienen tiempo de adaptarse \cite{book:zuniga2021}.
Las precipitaciones, además, no son uniformes, sino que aumentan en zonas donde ya son muy abundantes y disminuyen
en las regiones más secas \cite{book:zuniga2021}.

Ante esta situación se plantea la cuestión de cuáles de estos efectos podríamos ser capaces de predecir
y en qué proporción para, si bien no subsanar completamente, sí ser capaces de paliar en cierta medida
conociendo con anterioridad ciertos datos climatológicos que ayudarán a las diferentes tomas de decisiones de 
los viticultores.

De esta manera, siendo conocedores a corto o medio plazo la humedad y la temperatura superficial del suelo,
entre otras variables, podríamos adaptar los sistemas de riego para adecuarlos a las necesidades de los
cultivos y por otro lado, gestionar de forma adecuada los recursos hídricos nacionales (en este caso
de la zona hidrográfica del Duero) al mismo tiempo que protegemos a las vides de plagas y trastornos fisiológicos 
debidos a diferentes factores como la humectación o la temperatura del suelo.

En los últimos años las mejoras en los modelos predictivos basados en los datos, han hecho posible la aplicación
de diferentes soluciones en multitud de aplicaciones, ayudando a resolver distintos problemas relacionados con 
diversos campos, como la evolución futura de precios en el mercado bursátil, recomendación de productos, diagnósticos médicos, detección de fraudes, etc.

En este proyecto se propone el análisis de los datos de temperatura y humedad del suelo procedentes de 
sensores IoT (\textit{Internet of Things}) desplegados en un viñedo situado en la zona sur de la provincia de Burgos, para, de esta forma,
obtener modelos predictivos que nos permitan conocer el comportamiento de ciertas variables en función de otras
y anticipar su evolución futura a corto plazo.

Se emplearán diferentes modelos, siendo \textit{Long Short-Term Memory} 
(\textit{LSTM}) y \textit{Gated Recurrent Unit} (\textit{GRU}) las estructuras recurrentes estudiadas y 
\textit{Multi-Layer Perceptron} (\textit{MLP}) la contraparte convencional, realizando una comparación
entre las distintas soluciones para observar su comportamiento con los datos proporcionados.

\section{Estructura de la memoria}
La memoria del proyecto cuenta con la siguiente estructura:
\begin{itemize}
    \item Introducción: descripción del problema que el proyecto pretende resolver junto con la estructuración 
        de la memoria y los materiales adjuntos.
    \item Objetivos del proyecto: descripción de los objetivos a cumplir con el desarrollo del proyecto.
    \item Conceptos teóricos: explicación de los conceptos teóricos necesarios para la comprensión del proyecto, 
        el problema a abordar y la solución propuesta.
    \item Técnicas y herramientas: listado de técnicas y herramientas empleadas durante el desarrollo del proyecto.
    \item Aspectos relevantes del desarrollo: aspectos a destacar durante la realización del proyecto.
    \item Trabajos relacionados: trabajos relacionados con el problema que aborda el proyecto.
    \item Conclusiones y líneas de trabajo futuras: conclusiones obtenidas de la realización del proyecto y posibilidades
        de ampliarlo o de introducir mejoras.
\end{itemize}

Además de la memoria, se incluyen una serie de anexos:
\begin{itemize}
    \item Plan de proyecto software: planificación temporal y estudio de viabilidad del proyecto.
    \item Comprensión del negocio: especificación de los objetivos y requisitos que el proyecto pretende abordar.
    \item Comprensión de los datos: estudio del conjunto de datos inicial y su adecuación al problema.
    \item Preparación de los datos: proceso llevado a cabo para seleccionar los datos empleados en el modelado.
    \item Modelado: creación de los diferentes modelos.
    \item Evaluación: evaluación de los modelos obtenidos.
    \item Manual del programador: aspectos relevantes del código fuente del proyecto.
    \item Manual de usuario: guía de usuario para el manejo de la aplicación asociada al proyecto
\end{itemize}

\section{Materiales adjuntos}
El proyecto incluye el siguiente contenido:
\begin{itemize}
    \item Python Notebook con el pre-procesamiento de datos (denominado \textbf{data-preprocessing.ipynb}) y 
        las diferentes explicaciones.
    \item Python Notebook con la integración de los diferentes conjuntos de datos (\textbf{data-integration.ipynb})
    en un único fichero.
    \item Python Notebook con el modelado de la Red Neuronal Artificial empleada para la regresión
        de humedad y temperatura y las explicaciones correspondientes, denominado \textbf{ann-modeling.ipynb}.
    \item Conjunto de datos de los sensores desplegados en el viñedo.
    \item Conjunto de datos del pluviómetro desplegado en el viñedo.
    \item Scripts de Python para mostrar las gráficas de los datos (de sensores y pluviómetro) sin procesar.
    \item Scripts de Python para mostrar las gráficas de los datos (de sensores y pluviómetro) procesados
    \item Script de Linux para la instalación del entorno virtual.
    \item Script de Windows PowerShell para la instalación del entorno virtual.
    \item Script de Windows CMD para la instalación del entorno virtual.
\end{itemize}
Los recursos mencionados se encuentran disponibles en GitHub \cite{gabriel:2023}.

\capitulo{2}{Objetivos del proyecto}

El proyecto aborda diferentes objetivos. Podemos hacer una distinción entre 
objetivos generales, objetivos técnicos y objetivos personales.

\section{Objetivos generales}
\begin{itemize}
    \item Realizar un análisis de datos para ser capaces de predecir la temperatura
        y humedad del suelo en un viñedo.
    \item Facilitar la comprensión de los datos recogidos mediante representaciones
        gráficas.
    \item Buscar las correlaciones entre las diferentes medidas obtenidas y las 
        diferentes variables ambientales y/o de características del suelo.
    \item Encontrar modelos predictivos que permitan explicar el comportamiento
        de ciertas variables en función de las otras y anticipar la evolución futura
        a corto plazo.
\end{itemize}

\section{Objetivos técnicos}
\begin{itemize}    
    \item Realizar visualizaciones de los datos recogidos por los sensores IoT 
        (\textit{Internet of Things}) desplegados en el viñedo.
    \item Realizar un pre-procesamiento de datos mediante librerías de manipulación
        de datos de Python como Pandas.
    \item Modelar una Red Neuronal Artificial con Keras que permita predecir la humedad
        y temperatura del suelo gracias a los datos proporcionados.
    \item Emplear Git como sistema de control de versiones distribuido mediante
        la plataforma GitHub.
    \item Emplear ZenHub para la gestión de proyectos mediante las metodologías 
        ágiles.
    \item Aplicar en la medida de lo posible las metodologías ágiles mediante la 
        técnica Scrum aprendida durante el desarrollo del grado (ciertos aspectos
        como las ``\textit{daily}'' no pueden aplicarse debido al carácter del TFG).
\end{itemize}

\section{Objetivos personales}
\begin{itemize}
    \item Emplear los conocimientos y técnicas adquiridas durante el desarrollo de los diferentes
        cursos del Grado en Ingeniería Informática.
    \item Profundizar en el uso de la inteligencia artificial para la resolución de 
        problemas cotidianos que tendrá un reflejo en un sistema real (como es el
        campo del cultivo de la vid).
    \item Profundizar en la utilización de un lenguaje de programación tan versátil
        como es Python para el análisis de datos, empleando diferentes librerías
        nativas y de terceros, así como crear nuevos módulos.
    \item Profundizar en la librería de Tensorflow que se encuentra tan extendida
        en los diferentes ámbitos de la vida cotidiana en los que se emplean las redes 
        neuronales (mundo laboral e investigador). 
\end{itemize}
\capitulo{3}{Conceptos teóricos \cite{book:witten2017}}
El desarrollo del proyecto cuenta con diferentes fases siguiendo el proceso de \textit{descubrimiento de conocimiento
en bases de datos, KDD}, compuesto por la comprensión del negocio, la comprensión de los datos, la preparación de datos,
el modelado, la evaluación del modelo y el despliegue del producto software.

\imagen{data_mining_process.jpg}{\textit{Proceso KDD}. Extraido de \cite{book:witten2017}}{0.5}

En lo referente a la fase de comprensión del negocio se analizan los objetivos y requisitos del proyecto. 
En este caso quedan bien marcados en la descripción del proyecto y la introducción, así como en la sección de los objetivos.

En la etapa de comprensión de los datos se crea el conjunto de datos inical y se comprueba si este es adecuado,
de forma que si se determina que no lo es se deberán recopilar más.
En el caso que nos atañe como estamos sujetos a los plazos del curso la recopilación de más datos puede complicarse.

En la preparación de los datos se realiza el pre-procesamiento de estos, de forma que puedan ser empleados en el modelado.

En la fase de modelado se crean los modelos, lo que irá de la mano con la fase de preparación, puesto que algunas herramientas de 
pre-procesamiento incluyen un modelo interno de los datos para transformarlos.

En la fase de la evaluación se estima el rendimiento del modelo, reconsiderándose en su caso los objetivos, de forma que si los modelos son poco efectivos
se vuelve a la primera fase.

No se contempla realizar un despliegue de los modelos resultantes en ninguna aplicación, puesto que se sale del alcance del proyecto.

En esta sección se presentarán los conceptos teóricos relevantes en cada etapa para facilitar su comprensión.

\section{Pre-procesamiento de datos}
En el pre-procesamiento de datos se pretende realizar la integración y limpieza de estos, de forma que disminuyan los posibles problemas
de calidad que puedan surgir en los diferentes sistemas de información.

Como norma general el \textbf{proceso de integración} debe ser realizado durante la fase de recopilación de los datos.
La \textbf{limpieza} permite la detección y corrección de los problemas no resueltos en la fase anterior como los valores anómalos (\textit{Outliers}) o faltantes 
\cite{book:hernandez2004}.  

Tras la integración de los datos de las diferentes fuentes (e.g. bases de datos), se pueden realizar un resumen de características de atributos, en la que
se mostrarán las características generales de estos como medias, mínimos, máximos y valores posibles. 

En esta tabla podríamos obtener información trascendental para proceso de análisis de datos, sobre todo para atributos categóricos.
En el caso de atributos numéricos un mecanismo visual que es especialmente útil es la gráfica de dispersión \cite{book:hernandez2004}, que es la técnica de visualización
de datos que se empleará mayoritariamente durante el desarrollo del proyecto.

En el conjuto de datos podemos encontrar \textbf{valores faltantes}, perdidos o ausentes que pueden ser reemplazados por diferentes razones. Una de ellas es que
el modelo que empleemos puede no tratar bien estos valores o que utilice un mecanismo de tratamiento que no sea adecuado.
Un problema asociado a su detección es que estos no estén representados como nulos, lo que puede introducir sesgo en el conocimiento
extraido \cite{book:hernandez2004}.

Ante esta situación se puede actuar de diferentes maneras \cite{book:hernandez2004}:
\begin{itemize}
    \item Ignorarlos (ciertos algoritmos son tolerantes a los valores faltantes).
    \item Eliminar la columna que contiene valores faltantes.
    \item Filtrar las filas: eliminar las filas afectadas, lo que introduce cierto sesgo en muchas ocasiones.
    \item Reemplazar el valor por otro que preserve la media y la varianza del conjunto de datos en caso de atributos numéricos y la moda en atributos
        catregóricos.
        Una forma de reemplazar los valores faltantes es la imputación de datos perdidos, que consiste en predecirlos a partir de otros ejemplos. Existen también
        algoritmos que se emplean tradicionalmente para este fin.
    \item Segmentar: se segmentan las filas por los valores disponibles y se obtiene un modelo por cada uno de los segmentos y se combinan.
    \item Modificar la política de calidad de datos y esperar a que los faltantes estén disponibles.
\end{itemize}

Además de las situaciones anteriores, es posible que el conjunto de datos cuente con \textbf{valores erróneos} que 
se deben detectar y tratar.
La detección de estos campos puede realizarse de diferentes maneras.

En el caso que nos atañe se deben buscar los \textbf{valores extremos}, que no significa que sean erróneos, sino que 
estadísticamente se clasifican como anómalos, aunque representen un estado genuino de la realidad.
Con todo y con eso, estos valores pueden suponer un problema para algunos métodos que se basan en el ajuste de pesos
como las redes neuronales.
En otras ocasiones pueden haber datos erróneos que caen en la normalidad, por lo que no pueden ser detectados \cite{book:hernandez2004}.

La falta de detección de estos valores puede resultar en problemas si posteriormente se normalizan los datos, puesto
que la mayoria de datos estarían en un rango muy pequeño, lo que puede producir una falta precisión en algunos modelos \cite{book:hernandez2004}.

El tratamiento de los datos erróneos o anómalos pueden ser tratados de forma similar a los faltantes \cite{book:hernandez2004}:
\begin{itemize}
    \item Ignorarlos (ciertos algoritmos son robustos a datos anómalos).
    \item Eliminar la columna que contiene los datos anómalos, por ejemplo si esta situación se produce
        continuamente en el atributo (es preferible reemplazarla por otra columna con valores discretos
        estableciendo la corrección o no del valor).
    \item Filtrar las filas: eliminar las filas afectadas, lo que puede introducir cierto sesgo.
    \item Reemplazar el valor por nulo, por los máximos o mínimos del atributo o por
        las medias (se debe tener en cuenta que los modelos sean capaces de procesar los valores
        nulos, ya que en caso contrario se deberá hacer frente a un nuevo problema en el conjunto
        de datos).
    \item Discretizar: transformar un valor continuo en uno discreto.
\end{itemize}
Los atributos con valores erróneos serán mas graves cuando este sea empleado como clase o valor de salida de la predicción \cite{book:hernandez2004},
puesto que afectarán directamente al cálculo del rendimiento del modelo. 

\section{Modelado}

Para la creación de los modelos emplearemos diferentes tipos de \textit{Redes Neuronales Artificiales}, como 
el perceptrón multicapa (\textit{MLP}), \textit{Gated Recurrent Unit} y \textit{Long Short-Term Memory}.

\subsection{MLP}
El perceptrón multicapa además de la capa de entrada y salida contiene varias capas intermedias
llamadas ocultas (\textit{hidden layers}), denominadas de este modo puesto que los cálculos
internos están ocultos al usuario \cite{book:aggarwal2018}.

En este tipo de redes las señales se propagan inicialmente hacia adelante para obtener un resultado, 
de forma que las neuronas de una capa se encuentran interconectadas con las de la capa
siguiente y así sucesivamente en todas ellas.

Cada una de estas conexiones cuenta con un peso, que se inicializa generalmente de forma aleatoria. Estos pesos son
empleados para el cálculo de las señales de activación de cada neurona, que se emplearán en la función de activación:

\begin{equation}
    \begin{aligned}
    & \bar{h}_1=\Phi\left(W_1^T \bar{x}\right) \\
    & \bar{h}_{p+1}=\Phi\left(W_{p+1}^T \bar{h}_p\right) \quad \forall p \in\{1 \ldots k-1\}\\
    & \bar{o}=\Phi\left(W_{k+1}^T \bar{h}_k\right)
    \end{aligned}
\end{equation}

Donde \begin{math}\bar{h}_1\end{math} será el resultado de aplicar la función de activación en la primera capa oculta (la capa de entrada
actua de intermedio entre las entradas y la \textit{RNA}), \begin{math}\bar{h}_{p+1}\end{math} serán los resultados de la función de activación
en la capa p-ésima y \begin{math}\bar{o}\end{math} en la capa de salida.

\imagen{mlp.png}{Perceptrón Multicapa. Extraido de \cite{book:aggarwal2018}}{0.5}

Existen diversas funciones de activación entre las que destacan las sigmoides por ser aplicadas de forma 
elemental \cite{book:aggarwal2018}.

Tras la propagación de las señales de activación hacia adelante, se calculará el gradiente del error empleando el algoritmo
de \textit{Backpropagation}, que utilizará la regla de la cadena del cálculo diferencial para realizar este propósito, de forma 
que los gradientes calculados se emplearán posteriormente para actualizar los pesos de las conexiones sinápticas.

\subsection{Redes Neuronales Recurrentes (\textit{RNN})}
Las redes neuronales convencionales (véase por ejemplo los \textit{MLP}, las redes neuronales de base radial, etc.) están diseñadas para
datos multidimensionales cuyos atributos son independientes de otros.
Sin embargo, existen datos que contienen dependencias de atributos previos como los datos de series temporales, 
biológicos o cadenas de texto \cite{book:aggarwal2018}.

Las redes neuronales recurrentes pueden ser empleadas en este tipo de datos y además cuentan con ventajas con respecto a las redes convencionales.
A diferencia de estas recibe y procesa las entradas en el orden de llegada y las trata de igual manera que los instantes anteriores.
Por otro lado, se considera unidad básica a las estructuras denominadas celdas.

Además, en luagar de un número variable de entradas, contienen un número variable de celdas, y cada una de estas se corresponde con un 
instante de tiempo, lo que permite a cada una de estas interactuar con otras más cercanas a la salida \cite{book:aggarwal2018}.

Cada una de las celdas de este tipo de modelos realiza operaciones con dos entradas y dos salidas que se calculará en cada instante de tiempo \textit{t} como:
\begin{equation}
    \begin{aligned}
    f: \quad \mathbb{R}^{m+n} & \rightarrow \mathbb{R}^{n+l} \\
    \left(x_t, s_{t-1}\right) & \mapsto\left(y_t, s_t\right)
    \end{aligned}
\end{equation}

Donde \(x_t\) es el vector de entrada de dimensión \textit{m} a la celda en el instante \textit{t}, \(y_{t}\) es el vector de salida de 
dimensión \textit{n} y \(s_{t}\) es el vector de dimensión \textit{l} que representa el estado actual de la red.
Estas celdas compartirán los parámetros internos, por lo que la diferencia entre un instante u otro será las entradas a la celda.

\imagen{rnn.png}{Red Neuronal Recurrente.}{0.25}

En la Figura \ref{fig:rnn.png} puede observarse la representación de una celda de una red recurrente. En
esta puede verse cómo dada una secuencia en el instante \textit{t}, se produce una salida 
\textit{y} que generará una secuencia en el tiempo compartiendo parámetros (pesos sinápticos) como puede apreciarse
en la Figura \ref{fig:rnn_tiempo.png}.

\imagen{rnn_tiempo.png}{Representación de la RNN en el tiempo.}{0.55}

La salida del modelo en un instante de tiempo se calculará, de esta forma cómo:
\begin{equation}
    y^t = f(W^{yx}x^t + U^{yy} y^{t-1} + b) 
\end{equation}

Donde \(W^{yx}\) son los pesos sinápticos en el instante actual en la celda, mientras que \(U^{yy}\) se
trata de los pesos de la conexión recurrente, es decir, depende de la sinapsis del instante anterior.
La primera matriz de pesos, de esta forma, tendrá un tamaño de \(|x| \cdot |y|\), mientras que la segunda
será de tamaño \(|y| \cdot  |y|\). De manera similar \(|b| = |y|\), tratándose este del vector de sesgo.

A diferencia del modelo expuesto con anterioridad, las redes neuronales recurrentes no pueden aplicar el algoritmo de retropropagación (\textit{Backpropagation})
directamente, puesto que en este caso las celdas de instantes posteriores se encontrarán directamente conectadas con las anteriores.
En este caso se emplea una adaptación conocida como \textit{Backpropagation through time}, que consiste en ``desenrollar'' la red recurrente, convirtiéndola
en una red \textit{feed-forward} y posteriormente aplicar el algoritmo de retropropagación, teniendo en cuenta que dependiendo los instantes 
que se consideren como salidas de la red el gradiente puede afectar a un mismo parámetro de la red de forma repetida \cite{book:rue2019}.

\subsubsection{LSTM}
Las redes neuronales recurrentes tradicionales pueden encontrar problemas en capturar y recordar información
a largo plazo debido al problema denominado \textit{desvanecimiento o explosión del gradiente}, que ocurre
cuando los gradientes disminuyen o aumentan excesivamente a medida que se retropropagan a través
del tiempo, produciendo problemas en el aprendizaje a largo plazo.

Para abordar estos problemas, \textit{LSTM} introduce una estructura de memoria adicional denominada
celda de memoria que permite al modelo recordar información durante largos periodos de tiempo. Estos 
modelos emplearán para ello puertas que controlan el flujo de información dentro de la red.

Cada celda cuenta con tres puertas: la puerta de entrada \textit{Input Gate}, la puerta de salida 
(\textit{Output Gate}) y la puerta de olvido (\textit{Forget Gate}). La primera decide
cuánta información nueva debe debe agregar a la memoria, la segunda cuánta transmitir como salida y 
la tercera cuánta información antigua se debe de descartar.

En las celdas del modelo se calculan 3 funciones:
\begin{equation}
    f_t = \sigma (W^{f} x_t + U^{f} y_{t-1} + b_f)
\end{equation}
\begin{equation}
    i_t = \sigma (W^{i} x_t + U^{i} y_{t-1} + b_i)
\end{equation}
\begin{equation}    
    o_t = \sigma (W^{o} x_t + U^{o} y_{t-1} + b_o)
\end{equation}

La función \(f_t\) se corresponde con la puerta de olvido (\textit{Forget Gate}), mientras que 
\(i_t\) con la compuerta de entrada (\textit{Input Gate}) y \(o_t\) con la de salida (\textit{Output
Gate})

De esta manera, cada una de estas empleará dos matrices de pesos diferentes (\(W\) y \(U\)), que se compartirán
en cada instante de tiempo.

Con los diferentes valores de las compuertas, posteriormente se calcula un candidato de la cantidad de información
a añadirse:
\begin{equation}
    \widetilde{c}_t = tanh (W^{c} x_t + U^{c} y_{t-1} + b_c) 
\end{equation}

Finalmente el estado de la red se actualiza de la siguiente forma:
\begin{equation}
    c_t = f_t \odot  c_{t-1} + i_{t} \odot \widetilde{c}_t 
\end{equation}

Una vez calculado el nuevo estado, se emplea la puerta de salida, junto con la tangente hiperbólica
del nuevo estado oculto de la celda para obtener la nueva salida de la celda:
\begin{equation}
    y_t =  o_t \odot tanh(c_t)
\end{equation}

En la Figura \ref{fig:lstm_cell_mem.png} puede observarse el diagrama con las operaciones mencionadas
para producir la salida de la celda en el modelo \textit{LSTM}
\imagen{lstm_cell_mem.png}{Celda de memoria de LSTM}{1}

\subsubsection{GRU}


\capitulo{4}{Técnicas y herramientas}

Esta parte de la memoria tiene como objetivo presentar las técnicas metodológicas y las herramientas de desarrollo que se han utilizado para llevar a cabo el proyecto. Si se han estudiado diferentes alternativas de metodologías, herramientas, bibliotecas se puede hacer un resumen de los aspectos más destacados de cada alternativa, incluyendo comparativas entre las distintas opciones y una justificación de las elecciones realizadas. 
No se pretende que este apartado se convierta en un capítulo de un libro dedicado a cada una de las alternativas, sino comentar los aspectos más destacados de cada opción, con un repaso somero a los fundamentos esenciales y referencias bibliográficas para que el lector pueda ampliar su conocimiento sobre el tema.



\capitulo{5}{Aspectos relevantes del desarrollo del proyecto}

En este apartado se describirán los aspectos más importantes surgidos a partir del desarrollo del proyecto,
estos incluyen las complicaciones encontradas, así como las decisiones tomadas para solucionar todos los
problemas surgidos y sus implicaciones.

La sección se orientará siguiendo una estructura parecida a la introducida en \textit{Conceptos teóricos}, 
es decir, desarrollando los aspectos relevantes según la fase del proyecto en la que se produjo.

\section{Inicio del Proyecto}
Una de las motivaciones que se tuvo en cuenta para el desarrollo del proyecto fue la cercanía del tema tratado
con el lugar que se desarrolló. Y es que la provincia de Burgos se encuentra en la denominación de
origen ``Ribera del Duero'', una de las regiones que cuenta con algunos de los mejores vinos del mundo \cite{misc:abc2021}.

Otro impulso que llevó a realizar el Trabajo Final de Grado con este tema fue, de igual forma, la importancia 
del vino en la familia, pues el sustento de parte de esta es gracias a la viticultura, además de la historia
e importancia de la cultura del vino en el todo el territorio español.

Tras contactar con los docentes, se proporcionó el conjunto de datos empleado en el proyecto. En estas conversaciones
iniciales se consideró que estos contaban con ciertas características que debían tenerse en cuenta: en primer lugar,
los presentaban valores faltantes debidos a la posible desconexión de los sensores de la red, lo que provocaba
que estos dejaran de enviar datos al servidor. Por otro lado, las labores del campo podían haber afectado a la calidad
de estos por la extracción de las sondas de los sensores de sus posiciones originales.

Además, el pluviómetro instalado en un orígen se encontraba anclado de forma débil, lo que propiciaba que el viento 
produjera lecturas erróneas en el movimiento del balancín empleado en el registro de las precipitaciones.

Tras estas reuniones se inició el proceso de la comprensión del negocio, es decir, se comenzó con la definición de los 
objetivos del proyecto y los requisitos. Simplificando lo mencionado en la introducción y los objetivos,
la finalidad era analizar los datos proporcionados (centrándose en la temperatura y humedad) realizando
las visualizaciones pertinentes y encontrar modelos predictivos que intenten anticipar la evolución a corto plazo
empleando diferentes técnicas como las Redes Neuronales Artificiales.

\section{Fase de preparación de los datos}
En la fase de preparación se determinó que los datos proporcionados contaban con diferentes situaciones. En algunos
de los sensores la calidad de estos era bastante adecuada a pesar de contar con valores faltantes, en otros, en cambio,
existían numerosos problemas.

Por un lado, contaban con valores extremos que debían tenerse en cuenta para que los modelos que se desarrollasen no 
encontraran problemas al generalizar, desviándose de ese propósito por el posible ruido introducido.
Para solucionar este problema se empleó el mecanismo de detección de \textit{outliers} mediante el rango intercuartílico,
sustituyendo los datos identificados como extremos con la mediana del día en el que se realizó la muestra.

Por otro lado, algunos de los sensores contaban con valores matemáticamente correctos (dentro de límites lógicos de las 
diferentes variables) pero que contextualmente eran incorrectos. Un ejemplo de esto eran las lecturas artificiales de 
distintas variables como la temperatura cuando el nivel de la batería de los sensores caía por debajo de 3,3V, lo que
probocaba un desbordamiento de la variable en el controlador correspondiente.

Ante esta situación se decidió eliminar aquellas muestras cuyo nivel de batería se encontrara por debajo del mencionado
y por encima de los 4,1V (el máximo). Sin embargo, tras realizar esta operación se observó que ciertas 
muestras continuaban siendo contextualmente incorrectas arrojando datos inverosímiles (bajadas de temperatura y/o 
humedad extremas en periodos de tiempo cortos) debidas en gran medida, como se confirmó posteriormente, 
a la extracción de las sondas de la situación original, por lo que se acordó con los docentes la eliminación de estos 
valores que continuaban siendo erróneos.

El proceso se realizó a mano teniendo en cuenta las situaciones localizadas con anterioridad y con las nuevas necesidades,
marcando en los diferentes ficheros la validez de la entrada (cabe destacar que la recuperación de estos datos no era viable 
por las condiciones en la que se encontraban, puesto que el lapso temporal era demasiado elevado).

Además de estas condiciones, algunos de los sensores contaban con muestras con discontinuidades inexplicables 
mediante los datos, por lo que se acordó el intento de la recuperación de estos. Para ello, se observaron los muestreos
de los sensores geográficamente más cercanos al afectado, de forma que se pudiera conocer de manera más 
precisa cúal de las series era la que contaba con la lecura correcta y acutar en consecuencia.

A pesar de que se haya podido introducir algún sesgo debido a la modificación de los valores plasmados en los muestreos, 
los datos reales serán más cercanos a los nuevos valores asignados que a las lecturas anteriores, es decir, aún si 
estos fueran incorrectos serían menos incorrectos que anteriormente.

Ante la premura de los plazos del proyecto, realizar de nuevo el muestreo de los datos suficientes para crear un modelo predictivo
eficiente no era viable, por lo que se tuvo que trabajar con la cantidad proporcionada.

\section{Fase de creación de modelos}
Tras realizar una investigación preliminar acerca de los modelos con mayor precisión 
en la regresión de series temporales, se llegó a la conclusión de emplear
redes neuronales recurrentes con memoria como primer objetivo de modelado, 
puesto que sobre datos meteorologicos contaban con un buen rendimiento.

Por ese motivo, se escogió utilizar las redes neuronales con memoria a corto y medio
plazo como son \textit{GRU} y \textit{LSTM}, empleando los mismos hiperparámetros y,
de esta manera, poder realizar una comparativa para tener una idea de cómo se 
ajustan estos modelos a las diferentes funciones objetivo.

Además, otra opción a emplear fue la tradicional \textit{MLP} para poder observar
cómo la linealidad de este modelo (no tiene en cuenta instantes anteriores) impide obtener unos 
resultados aceptables en comparación de sus contrapartes recurrentes.

El primer reto consistía en preparar los datos de entrada a la red neuronal y los 
correspondientes salidas de forma que estos fueran muestras homogéneas
sin saltos temporales que impidiearan a los modelos ajustarse adecuadamente.

En primer lugar, en un espacio de 10 minutos se podían encontrar 2 o 3 muestras según fuera 
el \textit{offset} del controlador de muestreos de los sensores.
Además, como en la fase de preparación de los datos se realizaron descartes, se generaron, de esta manera,
saltos temporales que podían afectar a la predicción de la tendencia.

Una solución fue recorrer el conjunto de datos del sensor con una ventana deslizante del tamaño
equivalente a las instancias de entrada (establecido por parámetro) y el dato de salida (se realizó una 
aproximación preliminar realizando una única predicción, que posteriormente se modificó para 
que los modelos fuesen capaces de realizar varias), de esta forma en ventanas en las que 
no habían datos suficientes para la entrada se descartaban.

\imagen{window.png}{Esquema del proceso de preparación de datos}{0.5}

Cabe destacar que inicialmente se emplearon todos los datos para entrenamiento y validación, no estableciendo
el conjunto de test, lo que fue modificado más tarde añadiendo el tercer conjunto, permitiendo reducir el sesgo en 
la evaluación del modelo, pues, a diferencia de las instancias de validación, se tratan
de datos que no han sido empleados por el modelo tanto para ser entrenado como validado.

Tras la preparación de los datos de entrada de la red se comenzó a realizar el modelo \textit{GRU}, de forma
que se debían establecer sus capas y las celdas en estas. Posteriormente se realizaron los modelos
\textit{LSTM} y \textit{MLP}.
Para agilizar la construcción y verificar su funcionamiento se estableció inicialmente 3 capas en todos ellos, con
36, 16 y 8 unidades sucesivamente.

A diferencia de los dos primeros, \textit{MLP} requería de una ligera modificación, puesto que no aceptaba secuencias, 
por lo que se estableció una neurona de entrada por cada valor en la serie de tiempo y de igual forma una 
neurona en la capa de salida por cada predicción. 

\section{Fase de evaluación de modelos}
En la fase de evaluación se estableció realizar pruebas sobre los modelos modificando ciertos hiperparámetros como 
el número de capas ocultas, el número de unidades (celdas o neuronas dependiendo del contexto) por capa y
ratio de aprendizaje.

De esta forma, los modelos empleados tendrán las especificaciones estándar en el caso de
\textit{GRU} y \textit{LSTM}, modificando la función de activación de las neuronas en las capas ocultas en 
\textit{MLP} por una función \textit{ReLU} (Figura \ref{fig:relu.png}), que aportará la no linealidad al modelo, lo que presumiblemente
mejoraría el rendimiento concreto con los datos empleados.

\imagen{relu.png}{Función ReLU}{0.75}

Se realizarán, entonces, combinaciones con diferentes valores:
\begin{itemize}
    \item Capas: 1, 2 y 3 capas.
    \item Tamañdo de capa: 8, 16 y 32 celdas/neuronas.
    \item Ratio de aprendizaje: 0.1, 0.01 y 0.001.
\end{itemize}

\capitulo{6}{Trabajos relacionados}

Este apartado sería parecido a un estado del arte de una tesis o tesina. En un trabajo final grado no parece obligada su presencia, aunque se puede dejar a juicio del tutor el incluir un pequeño resumen comentado de los trabajos y proyectos ya realizados en el campo del proyecto en curso. 

\capitulo{7}{Conclusiones y Líneas de trabajo futuras}

Todo proyecto debe incluir las conclusiones que se derivan de su desarrollo. Éstas pueden ser de diferente índole, dependiendo de la tipología del proyecto, pero normalmente van a estar presentes un conjunto de conclusiones relacionadas con los resultados del proyecto y un conjunto de conclusiones técnicas. 
Además, resulta muy útil realizar un informe crítico indicando cómo se puede mejorar el proyecto, o cómo se puede continuar trabajando en la línea del proyecto realizado. 



\bibliographystyle{plain}
\bibliography{bibliografia}

\end{document}
